\cotefile{87}
\scanpage{3}
% archive sheet number: 2

\uline{Géométrie des convexes}

$\ast$ Partie convexe d'un espace affine réel : définition, stabilité par intersection; partie convexe engendrée - description en termes de combinaisons barycentriques, en termes de simplexes, comparaison avec espace affine engendré.

\inserted{Enveloppe convexe \unsure{de|d'} \deleted{\unsure{une}} \deleted{\unsure{famille|réun.}}famille d'ens. convexes.}

Stabilité par images inv et directes par appl. affines, par somme.

$\ast$ Cas de \deleted{$\unsure{\mathbb R|R}$}\inserted{dim $1$} : parties convexes sont les \uline{intervalles}.

\unsure{Une|\deleted{Une}} partie \deleted{\unsure{convexe}}\inserted{$C$} d'un \unsure{espace} affine réel $E$ est convexe ssi son intersection avec tte droite est un intervalle de ladite.

Dimension d'un convexe.

$\ast$ Intérieur d'un convexe (engendrant $E$ affinement), cas dim $E$

\[
\text{a)}\quad \mathring C\ne\varnothing
\]

\[
\begin{aligned}
\text{b)}\quad &\text{Si }x\in\mathring{\unsure{C|c}},\ \text{alors}\quad
 \mathring C=\bigcap_{\substack{D\ \text{droite}\\\text{de }E\ \text{passant}\\\text{par }x}}\operatorname{int}_{D}(D\cap C)\\
&\hspace{4em}\dot C=\bigcap_D\operatorname{bnd}_D(D\cap C)\\
&\hspace{4em}\overline C=\bigcap_D\overline{D\cap C}
\end{aligned}
\]

\margin{\deleted{$\left\{\begin{aligned}&\text{\unsure{faux, mais}}\\&\text{\unsure{juste si}}\\&C\ \text{fermé}\end{aligned}\right.$}}

c) Si $C$, \inserted{$C'$ convexes} compacts\unsure{,|;} \inserted{engend$^{\mathrm{nt}}$ $E$, $x\in\mathring C$, $x'\in\unsure{\dot C'|\mathring C'}$,} $\exists$ homéom. unique $h:C\simeq C'$ avec $h(x)=x'$, $h$ étant \deleted{\unsure{linéaire}}\inserted{affine} sur tte demi-droite issue de $x$ et la transformant en demi-droite parallèle
